\documentclass[11pt,a4paper]{article}

\usepackage[utf8]{inputenc}
\usepackage[T1]{fontenc}
\usepackage[french]{babel}
\usepackage{geometry}
\usepackage{setspace}
\usepackage{hyperref}
\usepackage{enumitem}

\geometry{margin=2.5cm}
\setstretch{1.15}

\title{Synthèse de la réunion du projet fil rouge \\ 
Axes d'amélioration, compléments et réorientations}
\author{Benjamin Lepourtois, ChatGPT}
\date{09 janvier 2026}

\begin{document}
\maketitle

\section{Problème central identifié}

Le point principal soulevé par les différents interlocuteurs concerne un \textbf{manque de problématisation explicite} du sujet.

L’état de l’art et la présentation globale du projet sont perçus comme :
\begin{itemize}
    \item trop orientés \textit{solution} (LLM $\rightarrow$ Behavior Trees),
    \item insuffisamment ancrés dans une \textbf{question scientifique clairement formulée},
    \item trop synthétiques pour permettre de comprendre les verrous réels du problème.
\end{itemize}

En conséquence, il est difficile pour les encadrants d’identifier précisément :
\begin{itemize}
    \item la ou les questions de recherche adressées,
    \item les verrous scientifiques ciblés,
    \item la cohérence entre ces verrous et la méthodologie proposée.
\end{itemize}

\section{État de l’art : restructuration nécessaire}

\subsection{Définition formelle des Behavior Trees}

Il est jugé indispensable d’introduire une définition claire et formelle des Behavior Trees, même si cela peut sembler redondant pour un public expert.

Les éléments suivants doivent apparaître explicitement :
\begin{itemize}
    \item définition d’un Behavior Tree,
    \item types de nœuds (Sequence, Fallback, Parallel, etc.),
    \item distinction entre actions et conditions,
    \item syntaxe XML associée,
    \item lien entre XML et moteur d’exécution (ex. BehaviorTree.CPP en C++),
    \item notion de validité syntaxique et de validité sémantique.
\end{itemize}

Cette étape permet :
\begin{itemize}
    \item d’établir une base conceptuelle commune,
    \item de vérifier la maîtrise de l’objet d’étude,
    \item de préparer rigoureusement le lien avec la génération automatique par LLM.
\end{itemize}

\subsection{Clarification des notions de syntaxe, sémantique et hallucination}

Les termes suivants sont utilisés mais insuffisamment définis :
\begin{itemize}
    \item \textbf{Syntaxe} : respect du schéma XML, structure valide, nœuds autorisés.
    \item \textbf{Sémantique} : comportement réel induit sur le robot, adéquation avec l’intention, absence d’événements redoutés.
    \item \textbf{Hallucination} : génération de skills inexistants, d’actions incohérentes ou de structures non exécutables.
\end{itemize}

Une clarification explicite est attendue afin d’éviter toute ambiguïté conceptuelle.

\subsection{Repositionnement de l’état de l’art}

L’état de l’art est actuellement perçu comme trop générique et centré sur les LLM.

Il est recommandé de :
\begin{itemize}
    \item partir des objets manipulés : Behavior Trees, XML, structures formelles,
    \item élargir à la littérature sur la génération de structures contraintes (XML, code, graphes),
    \item revenir ensuite au cas spécifique des Behavior Trees.
\end{itemize}

L’objectif est de montrer que le sujet s’inscrit dans un cadre plus large que le seul couplage BT--LLM.

\section{Méthodologie : justification des choix}

\subsection{Faire émerger les questions avant les solutions}

Les choix méthodologiques (fine-tuning, RAG, prompt système, grammaire formelle) apparaissent trop tôt.

Il est attendu de :
\begin{enumerate}
    \item formuler explicitement les questions scientifiques (ex. validité syntaxique, intégration de règles implicites, gestion de l’ambiguïté),
    \item montrer ensuite comment chaque technique répond à une question précise.
\end{enumerate}

\subsection{Rôle exact du LLM}

Il est nécessaire d’insister clairement sur le fait que :
\begin{itemize}
    \item le LLM ne crée pas de nouvelles compétences,
    \item il sélectionne et orchestre des skills existants,
    \item il agit comme un traducteur entre intention utilisateur et structure formelle.
\end{itemize}

Ce point est central pour la compréhension globale de l’approche.

\subsection{Gestion de l’implicite utilisateur}

La problématique de l’implicite est jugée pertinente mais insuffisamment visible.

À expliciter :
\begin{itemize}
    \item distinction entre information explicite et contraintes implicites (sécurité, règles métier),
    \item rôle du RAG pour injecter des règles non exprimées,
    \item possibilité d’interactions de clarification avec l’utilisateur.
\end{itemize}

\section{Évaluation : re-priorisation recommandée}

\subsection{Limites de l’évaluation fonctionnelle}

L’évaluation fonctionnelle complète est identifiée comme extrêmement coûteuse :
\begin{itemize}
    \item explosion combinatoire,
    \item dépendance forte à la simulation,
    \item complexité temporelle et capteurs bruités,
    \item difficulté à faire émerger des événements redoutés réalistes.
\end{itemize}

Il est recommandé de :
\begin{itemize}
    \item prioriser la validation syntaxique et structurelle,
    \item limiter l’évaluation fonctionnelle à des cas simples et artificiels,
    \item éviter de promettre une validation fonctionnelle réaliste exhaustive.
\end{itemize}

\section{Améliorations de la présentation}

\subsection{Manque de contextualisation}

Les slides sont jugées trop synthétiques.

Il est conseillé d’ajouter :
\begin{itemize}
    \item schémas explicatifs,
    \item exemples concrets de Behavior Trees,
    \item exemples XML minimaux.
\end{itemize}

\subsection{Séparation plus nette des parties}

Les frontières entre les sections doivent être clarifiées :
\begin{itemize}
    \item état de l’art,
    \item problématique,
    \item méthodologie,
    \item preuve de concept,
    \item limites et perspectives.
\end{itemize}

\section{Points positifs reconnus}

Les éléments suivants ont été explicitement valorisés :
\begin{itemize}
    \item preuve de concept pertinente et prometteuse,
    \item bonne appropriation de ROS2, Gazebo et BehaviorTree.CPP,
    \item intuition solide sur l’usage du RAG, des prompts systèmes et des grammaires,
    \item organisation de projet jugée sérieuse.
\end{itemize}

Le message global des encadrants est le suivant :
\begin{quote}
    La direction est pertinente, mais le cadrage scientifique et pédagogique doit être renforcé.
\end{quote}

\section{Checklist synthétique}

\begin{enumerate}
    \item Reformuler une problématique centrale claire
    \item Refaire l’état de l’art en partant des Behavior Trees et du XML
    \item Définir précisément syntaxe, sémantique et hallucination
    \item Justifier chaque choix méthodologique par une question
    \item Réduire les ambitions d’évaluation fonctionnelle
    \item Ajouter des exemples concrets et visuels
\end{enumerate}

\end{document}
